% ============================================================
% setup/preamble.tex —— 宏包与全局设置
% 本文件由 main.tex 通过 \input 引入，不要单独编译。
% 原则：只放"依赖与全局开关"，不放具体格式与内容。
% ============================================================

% ---------- 数学与符号 ----------
\usepackage{amsmath, amssymb, amsthm}
\usepackage{mathtools}
\usepackage{bm}                 % 粗体数学符号

% ---------- 图表 ----------
\usepackage{graphicx}           % 插图
\graphicspath{{figures/}{tables/}}  % 图片/表格搜索路径
\usepackage{subcaption}         % 子图（subfigure）
\usepackage{booktabs}           % 三线表
\usepackage{multirow}           % 表格跨行
\usepackage{makecell}           % 单元格内换行
\usepackage{longtable}          % 跨页长表
\usepackage{diagbox}            % 斜线表头（按需）
\usepackage{array}

% ---------- 页面尺寸与页眉页脚 ----------
\usepackage{geometry}           % 纸张与页边距
\usepackage{fancyhdr}           % 页眉/页脚定制（中国地质大学格式）

% ---------- 算法与代码 ----------
\usepackage{algorithm}
\usepackage{algorithmicx}
\usepackage{algpseudocode}
\usepackage{listings}           % 代码 listings
\usepackage{xcolor}

\lstset{
  basicstyle   = \small\ttfamily,
  breaklines   = true,
  frame        = single,
  numbers      = left,
  numberstyle  = \tiny,
  keywordstyle = \color{blue!70!black},
  commentstyle = \color{green!50!black},
  stringstyle  = \color{orange!80!black},
  captionpos   = b
}

% ---------- 参考文献（GB/T 7714-2005 国标，见规范 2.6）----------
% 顺序编码制：序号以方括号标注，置于引用处右上角（上标）。
% 正文引用请用 \supercite{key}（上标）；gb7714-2005 同时支持 \cite（行内）。
\usepackage[
  backend  = biber,
  style   = gb7714-2005,
  gbalign = left,
  sorting = none
]{biblatex}
\addbibresource{bib/references.bib}

% ---------- 交叉引用与超链接 ----------
\usepackage{hyperref}           % 建议最后加载此类宏包
\hypersetup{
  colorlinks = true,
  linkcolor  = black,
  citecolor  = black,
  urlcolor   = black,
  pdfborder  = {0 0 0}
}
\usepackage{cleveref}           % 智能引用（在 hyperref 之后）
\usepackage{siunitx}            % 单位与数字排版（规范 2.7：GB 3100~3102）
  \sisetup{
    inter-unit-product = \ensuremath{{}\cdot{}},  % 单位间用点乘
    per-mode = symbol,                             % 用 / 表示"每"
    separate-uncertainty = true
  }
\usepackage{enumitem}           % 列表环境定制
